
\section{An algorithm for Large-Integer Multiplication}\doHeader{Large-Integer Multiplication}

\paragraph{Problem definition}
\begin{mylist}
\item \textbf{input:}
  $n$-bit integers $A$ and $B$,
  with binary representations $A=a_1a_2\ldots a_n$
  and $B=b_1b_2\ldots b_n$.
\item \textbf{output:} the product $C = A \, B$.
\end{mylist}

This problem is about multiplying integers that are too large to fit in a machine word.
For example, RSA encryption currently requires arithmetic on 2048-bit integers, or larger,
while current machine words are typically 64 bits.

\paragraph{Algorithm---first attempt.}
The standard grade-school algorithm takes time $\Theta(n^2)$.
We want something faster.
We start by considering the following recursive algorithm:
\begin{myframed}
\begin{steps}
  \item[] {\sf int-product-1}$(n, A=a_1a_2\ldots a_n, B=b_1 b_2 \ldots b_n)$
    \comment{($A$ and $B$ are $n$-bit integers; $n$ is a power of two)}
    \step if $n = 1$: return $a_{1} b_{1}$
    \comment{(base case)}

    \step split $A$ into $A_1 = a_1 \cdots a_{n/2}$ and $A_2 = a_{n/2+1} \cdots a_n$
    so that \(A = 2^{n/2} A_1 + A_2\).

    \step split $B$ into $B_1 = b_1 \cdots b_{n/2}$ and $B_2 = b_{n/2+1} \cdots b_n$
    so that \(B = 2^{n/2} B_1 + B_2\).

    \step return $2^{n} (A_1 \times B_1) + 2^{n/2}(A_1 \times B_2 + A_2 \times B_1) + (A_2\times B_2)$
  \end{steps}
\end{myframed}

The last step requires four multiplications of $n/2$-bit integers.
\emph{These are done recursively!}

The two multiplications by powers of 2 are done directly just by shifting the bits.

\paragraph{Correctness.}
Correctness of the algorithm follows by induction from the equation
\begin{align}
  A\times B & = (2^{n/2} A_1 + A_2)\times (2^{n/2} B_1 + B_2) \notag\\
            & = 2^{n} (A_1 \times B_1)
              + 2^{n/2}(A_1 \times B_2 + A_2 \times B_1)
              + (A_2\times B_2).
              \label{dc: kar 1}
\end{align}
Because of this equation, as long as the four sub-products are calculated correctly
(and the additions are done correctly)
the result returned will be correct.

\paragraph{Run time.}
Each addition can be done using the standard algorithm in $O(n)$ time.
And each multiplication by a power of two just shifts the bits,
so can be done in $O(n)$ time.
So the running time satisfies $T(n) = 4 T(n/2) + n$.
The recursion-tree approach gives
\begin{align*}
  T(n)
  & {} = \sum_{i=0}^{\text{\#levels}} (\text{\# subproblems in level } i)
    \times (\text{work per subproblem in level } i) \\
  & {} = \sum_{i=0}^ {\log_2 n} (4^i) (n/2^i)
  = n \sum_{i=0}^{\log_2 n} 2^i
  = n\, \Theta(2^{\log_2 n})
  = \Theta(n^2).
\end{align*}

This isn't any faster than the standard algorithm!  Can we do better?

\paragraph{Karatsuba's algorithm.}
The middle term in Step 4, namely $A_1\times B_2 + A_2 \times B_1$, satisfies
\begin{equation}\label{dc: kar 2}
  A_1\times B_2 + A_2 \times B_1
  = (A_1 + A_2)\times (B_1 + B_2) - A_1\times B_1 - A_2 \times B_2.
\end{equation}
The terms $A_1\times B_1$ and $A_2 \times B_2$ are computed anyway
for Step 4, so by computing the right-hand side with this equation,
we replace two recursive calls by \emph{one}.
This gives us Karatsuba's algorithm:

\begin{myframed}
\begin{steps}
  \item[] {\sf karatsuba}$(n, A=a_1a_2\ldots a_n, B=b_1 b_2 \ldots b_n)$
    \comment{($A$ and $B$ are $n$-bit integers; $n$ is a power of two)}
    \step if $n = 1$: return $a_{1} b_{1}$
    \comment{(base case)}

    \step split $A$ into $A_1 = a_1 \cdots a_{n/2}$ and $A_2 = a_{n/2+1} \cdots a_n$
    so that \(A = 2^{n/2} A_1 + A_2\).

    \step split $B$ into $B_1 = b_1 \cdots b_{n/2}$ and $B_2 = b_{n/2+1} \cdots b_n$
    so that \(B = 2^{n/2} B_1 + B_2\).

    \step let $M_{1} = \textsf{karatsuba}(n/2, A_1, B_1)$
    \comment{($M_{1} = A_1 \times B_1$)}
    
    \step let $M_{2} = \textsf{karatsuba}(n/2, A_2, B_2)$
    \comment{($M_{2} = A_2 \times B_2$)}

    \step let $M_{3} = \textsf{karatsuba}(n/2, A_1 + A_2, B_1 + B_2)$
    \comment{($M_{3} = (A_1 + A_2) \times (B_1 + B_2)$)}
    
    \step return $2^{n}\, M_1 + 2^{n/2}(M_3 - M_1 - M_2) + M_2$
  \end{steps}
\end{myframed}

\paragraph{Correctness.}
Correctness follows by induction from Equations~\eqref{dc: kar 1} and~\eqref{dc: kar 2}.

\paragraph{Run time.}
Recalling that each addition or multiplication by a power of two
can be done in $O(n)$ time,
the running time now satisfies $T(n) = \mathbf{3}\, T(n/2) + n$,
which implies
\[
  T(n) = \sum_{i=0}^ {\log_2 n} (3^i) (n/2^i)
  = n \sum_{i=0}^{\log_2 n} (3/2)^i
  = n\, \Theta((3/2)^{\log_2 n})
  = \Theta(n^{\log_2 3})
  = \Theta(n^{1.58\ldots}). 
\]
(Above we use $x^{\log y} = y^{\log x}$ and $\log_2 (3/2) = \log_2(3) - 1$.) 
We have the following result:
\begin{theorem}[Karatsuba]
  There is an $O(n^{\log_2 3})$-time algorithm
  for multiplying $n$-bit integers.
\end{theorem}

Note: A recent (2019) algorithm achieves run time $O(n\log n)$.


\subsection{Exercise}

\exercise\emph{(Karatsuba's variation)}
The sums $A_1+A_2$ and $B_1+B_2$ may have $n/2+1$ bits due to carry.
This complicates the algorithm and the recurrence relation.
Professor Bo Zo proposes the following variant to fix this.  Does it work?
\begin{myframed}
\begin{steps}
  \item[] {\sf int-product-2}$(n, A=a_1a_2\ldots a_n, B=b_1 b_2 \ldots b_n)$
    \comment{($A$ and $B$ are $n$-bit integers; $n$ is a power of two)}
    \step if $n = 1$: return $a_{1} b_{1}$
    \comment{(base case)}

    \step split $A$ into $A_1 = a_1 \cdots a_{n/2}$ and $A_2 = a_{n/2+1} \cdots a_n$
    so that \(A = 2^{n/2} A_1 + A_2\).

    \step split $B$ into $B_1 = b_1 \cdots b_{n/2}$ and $B_2 = b_{n/2+1} \cdots b_n$
    so that \(B = 2^{n/2} B_1 + B_2\).

    \step let $M_{1} = \textsf{int-product-2}(n/2, A_1, B_1)$
    \comment{($M_{1} = A_1 \times B_1$)}
    
    \step let $M_{2} = \textsf{int-product-2}(n/2, A_2, B_2)$
    \comment{($M_{2} = A_2 \times B_2$)}

    \step let $M_{3} = \textsf{int-product-2}(n/2, |A_2 - A_1|, |B_2 - B_1|)$
    \comment{\color{red} $\leftarrow$ CHANGED THIS LINE AND BELOW!}
    
    \step let $s_A = \textsf{sign}(A_2 - A_1)$ and $s_B = \textsf{sign}(B_2 - B_1)$

    \step return $2^{n}\, M_1 + 2^{n/2}(M_1 + M_2 - s_A s_B M_3) + M_2$
  \end{steps}
\end{myframed}



\subsection{External sources on \prob{Large-Integer Multiplication}}

\begin{mylist}[-]
\item \DPV~Chapter 2.1, \KT Chapter 5.5

\item
  \JE Chapter 1, Section 1.9 (has information about the 2019 $O(n\log n)$-time algorithm)
  
\item An MIT lecture video
  
\URLsmall{https://ocw.mit.edu/courses/electrical-engineering-and-computer-science/6-006-introduction-to-algorithms-fall-2011/lecture-videos/lecture-11-integer-arithmetic-karatsuba-multiplication/}

\end{mylist}


\newpage


\section{An algorithm for Matrix Multiplication}\doHeader{Matrix Multiplication}
Next we describe a similar algorithm for Matrix Multiplication.

\paragraph{Problem definition}
\begin{mylist}
\item \textbf{input:}
  $n\times n$ matrices $A$ and $B$.
\item \textbf{output:} the matrix product $C = A \times B$.
  (So $c_{ij} = \sum_{k=1}^n a_{ik} b_{kj}$.)
\end{mylist}

We assume throughout that we are working with 64-bit integers,
so multiplication and addition of two numbers can be done in constant time.
We also assume $n$ is a power of 2
(we can reduce the general case to this one).

\paragraph{Algorithm---first attempt.}
The obvious algorithm takes time $O(n^3)$.
We want time $o(n^3)$.

To warm up, consider the following recursive algorithm.
Let $A_{11}$, $A_{12}$, $A_{21}$, $A_{22}$,
$B_{11}$, $B_{12}$, $B_{21}$, $B_{22}$,
$C_{11}$, $C_{12}$, $C_{21}$ and $C_{22}$
be $\frac{n}{2} \times \frac{n}{2}$
submatrices such that
\[
A = \left[ \begin{array}{ll} A_{11} & A_{12} \\A_{21} & A_{22}\end{array}\right],~{}
B = \left[ \begin{array}{ll} B_{11} & B_{12} \\B_{21} & B_{22}\end{array}\right],~{}
\text{ and }
C = \left[ \begin{array}{ll} C_{11} & C_{12} \\C_{21} & C_{22}\end{array}\right].
\]
We are given $A$ and $B$, and want to compute $C$.
The following matrix equation holds for each of the four pairs $i,j\in\{1,2\}$:
\begin{equation}\label{dc: matrix recurrence 1}
  C_{ij} = A_{i1} \times B_{1j} + A_{i2} \times B_{2j}.
\end{equation}
This gives the following recursive algorithm:

\begin{myframed}
\begin{steps}
  \item[] {\sf matrix-product-1}$(n, A, B)$
    \comment{($A$ and $B$ are $n\times n$ matrices; $n$ is a power of two)}
    \step if $n = 1$: return $[a_{11} b_{11}]$
    \comment{(base case)}

    \step for each of the four pairs $i,j\in \{1,2\}$,
    compute $A_{ij}$ and $B_{ij}$ as defined above
    \comment{(time $O(n^2)$)}

    \step for each of the four pairs $i,j\in \{1,2\}$:
    \begin{block}
      \step let $C_{ij}
      = \textsf{matrix-product-1}(n/2, A_{i1}, B_{1j})
      + \textsf{matrix-product-1}(n/2, A_{i2}, B_{2j})
      $
    \end{block}

    \step return $C = \left[ \begin{array}{ll} C_{11} & C_{12} \\C_{21} & C_{22}\end{array}\right]$
  \end{steps}
\end{myframed}

\paragraph{Correctness.}
Correctness of the algorithm follows
from Equation~\eqref{dc: matrix recurrence 1}, by induction.

\paragraph{Running time.}
To compute the product of two $n\times n$ matrices,
the algorithm makes eight recursive calls on $(n/2) \times (n/2)$ matrices.
The time spent outside of the recursive calls is $O(n^2)$.
So the running time $T(n)$ satisfies
$T(n) = 8\,T(n/2) + O(n^2)$.
Using the recursion-tree method, the total time is proportional to
\begin{align*}
  T(n) & = \sum_{i=0}^{\log_2 n} (8^i)\times (n/2^i)^2 
    = n^2 \sum_{i=0}^{\log_2 n} 2^i
  = n^2\, \Theta(2^{\log_2 n})
  = \Theta(n^3).
\end{align*}
%  (We use above that the sum is geometric,
% so its value is proportional to its largest term.
% Then we use $2^{\log_2 n} = n$.)
No better!

\paragraph{Strassen's Algorithm.}
To improve the running time, we'll reduce the number of recursive calls
from eight to seven.
To do this, replace Steps 3 and 4.1 with the following steps.
First, compute the following seven $\frac{n}{2} \times \frac{n}{2}$ matrices,
using seven recursive matrix multiplications:
\begin{align*}
  {M}_{1} & = ({A}_{11} + {A}_{22}) \times ({B}_{11} + {B}_{22})
            & {M}_{2} & = ({A}_{21} + {A}_{22}) \times {B}_{11} \\
  {M}_{3} & = {A}_{11} \times ({B}_{12} - {B}_{22}) 
            & {M}_{4} & = {A}_{22} \times ({B}_{21} - {B}_{11}) \\
  {M}_{5} & = ({A}_{11} + {A}_{12}) \times {B}_{22} 
            & {M}_{6} & = ({A}_{21} - {A}_{11}) \times  ({B}_{11} + {B}_{12}) \\
  {M}_{7} & = ({A}_{12} - {A}_{22}) \times ({B}_{21} + {B}_{22})
\end{align*}  
Strassen figured out that those seven matrices satisfy the following equations:
\begin{lemma}[Strassen]
  \begin{align*}
    {C}_{11} & = {M}_{1} + {M}_{4} - {M}_{5} + {M}_{7}
    & {C}_{12} & = {M}_{3} + {M}_{5} \\
    {C}_{21} & = {M}_{2} + {M}_{4}
    & {C}_{22} & = {M}_{1} - {M}_{2} + {M}_{3} + {M}_{6}
  \end{align*}
\end{lemma}
We don't give the proof, but you can check each equation if you like.

Now we modify the Algorithm in Step 4.1 so it computes each $C_{ij}$ using the four equations above
in time $O(n^2)$.  This gives Strassen's algorithm.

\paragraph{Correctness.}
Correctness of the algorithm follows by the lemma (and induction on $n$).

\paragraph{Run time.}
The time now $T(n)$ now satisfies
$T(n) = \mathbf{7}\,T(n/2) + O(n^2)$,
so the total time is proportional to
\[
  \sum_{i=0}^{\log_2 n} (7^i)\times (n/2^i)^2 
    = n^2 \sum_{i=0}^{\log_2 n} (7/4)^i
  = n^2\, \Theta((7/4)^{\log_2 n})
  = \Theta(n^{\log_2 7}) = \Theta(n^{2.80\ldots}).
\]
(Above we use $x^{\log y} = y^{\log x}$ and $\log_2 (7/4) = \log_2(7) - 2$.)
We have the following result:
\begin{theorem}[Strassen]
  There is an $O(n^{\log_2 7})$-time algorithm
  for multiplying $n\times n$ matrices.
\end{theorem}

\subsection{Exercise}

\exercise\emph{(Strassen variation)}
The standard way of multiply two $3\times 3$ matrices uses 27 multiplications.
Professor Bo Zo found a way to do it using 24
multiplications and a constant number of additions
(without assuming that multiplication is commutative).
He uses this observation as the basis for a variant of Strassen's algorithm
% [review 2026-09-27] fix: a 3x3 block scheme with 24 block products makes 24 recursive calls, not 80 (80 was inconsistent with "24 multiplications" and n/3 x n/3 blocks); the intended recurrence is T(n) = 24 T(n/3) + O(n^2)
that makes 24 recursive calls, each multiplying two matrices of size $n/3 \times n/3$.
Does he get a faster algorithm?
Show your work!

\subsection{External sources on fast matrix multiplication}

\begin{mylist}[-]

\item \CLRS Chapter 4.2 (includes exercises), and an MIT lecture video:
  
\URLsmall{https://ocw.mit.edu/courses/electrical-engineering-and-computer-science/6-046j-introduction-to-algorithms-sma-5503-fall-2005/video-lectures/lecture-3-divide-and-conquer-strassen-fibonacci-polynomial-multiplication/}

\end{mylist}

% \newcommand{\URL}[1]{\par-- \parbox[t]{0.95\textwidth}{\small\url{#1}} \medskip\par}


%%% Local Variables:
%%% mode: latex
%%% TeX-master: "divide_and_conquer"
%%% End:
